
\section{Introduction} \label{into}

This capstone project focuses on markerless 3D human pose tracking. The estimation and tracking of 3D human poses is an important research area in computer vision, driven by its broad range of immediate applications across multiple domains. For example, in robotics and human-robot interactions it can find application in industrial environments and collaborative workspaces where robots must accurately perceive human poses and predict movements to maintain safe operational boundaries and avoid collisions with dangerous machinery \parencite{svenstrup2009pose,dondrup2015real,koppula2015anticipating}. Similarly, for crowd control tracking of multi-person trajectories in public spaces facilitates anomalous behavior detection, crowd density planning, and automated surveillance \parencite{sultani2018real,luo2021multiple}. Other impactful applications can be found in healthcare, such as gait analysis, automated fall detection, and remote patient monitoring \parencite{han2025comprehensive,gaya2024deep}. Finally, human pose estimation and tracking can find applications also in the entertainment domain, such as the use of real-time motion capture for virtual avatars or natural body-based control interfaces for gaming \parencite{mehta2020xnect}.

Despite the maturity of real-time 2D human pose estimation, transferring these methods to a robust and real-time capable 3D setting remains a challenging task due to issues such as depth ambiguities, severe self-occlusions, and the high perspective variance of human motion \parencite{zheng2023deep,nogueira2025markerless}. The primary objective of this project is to design, implement, and evaluate a computer vision architecture capable of tracking the skeletal movement of multiple people from 2D pose estimations extracted from multiple views by 2D human pose estimation models. The system should be able to fuse the information gained from the 2D video frames of multiple camera viewpoints into coherent and temporally consistent 3D pose trajectories.

For the evaluation of the result, this project will use a subset of the CMU Panoptic dataset \parencite{Joo_2015_ICCV,Joo_2017_TPAMI}, which contains videos captured from multiple calibrated cameras together with ground truth pose and tracking information for each frame. Using this evaluation set this project further aims to investigate the impact of model capacity as well as the viewpoint count on both the achieved tracking accuracy and the runtime efficiency of the system. For qualitative analysis the project also makes use of the SALSA dataset \parencite{alameda2015salsa}. This dataset does not contain ground truth pose annotations, but represents a more difficult application scenario.

The remainder of this report is structured as follows: \Cref{background} reviews the relevant theoretical background underlying the project. \Cref{methodology} gives an overview of the developed tracking system as well as detailed explanations of each component. \Cref{experiments} describes the used dataset and the evaluation strategy including evaluation metrics. \Cref{results} presents the quantitative and qualitative evaluation results. Finally, \Cref{conclusion} concludes the report with a summary of the findings and outlines limitations, highlighting potential areas for further improvements.
